\scanpage{375}
\mspage{291}
% archive sheet number: 374

% slanted
\margin{le \unsure{rejeter|reporter} après 4.13.23}
\uline{Proposition} 4.13.7. Soit $M$ une catégorie accessible,
$(\unsure{\underline S|S}_\alpha)_{\alpha\in A}$ une \deleted{ensemble}\inserted{famille} de parties de $\operatorname{Ob}M$,
indexée par un petit ensemble $A$, $(\underline S_\alpha)_{\alpha\in A}$
la famille des sous-catégories accessibles dans $M$
correspondantes. Soit $S=\bigcap_{\alpha\in A}S_\alpha$, $\underline S=\bigcap_{\alpha\in A}\underline S_\alpha$, de
sorte que $\underline S$ est la sous-catégorie pleine de
$M$ définie par $S\subset\operatorname{Ob}M$. Ceci posé, si les
$S_\alpha$ sont accessibles, il en est de même
de $S$ -- ou encore, si les $\underline S_\alpha$ sont accessibles
dans $M$, il en est de $=$ de $\underline S$.

Cela résulte du résultat plus précis :

\uline{Corollaire} 4.13.8. Avec les notations \deleted{précéd. et \unsure{de}}\inserted{de 4.13.7 et 4.13.6,}
soit, pour tt $\alpha\in A$, $\pi'_\alpha$ un cardinal \inserted{$\geq\pi_1$} tel que
$\underline S_\alpha$ soit $\pi'_\alpha$-accessible.
% slanted
\margin{Cf reformulation plus \unsure{naturelle|rationnelle} et plus courante \unsure{en} pages 302, 302 bis \unsure{et|à} \unsure{ter}.}

\deleted{Soit \deleted{$\unsure{\pi'_0=\operatorname{Sup}(\operatorname{card}A,\operatorname{Sup}_\alpha 2^{({\pi'_\alpha}^{\pi_0})})}$}
\deleted{$\unsure{\pi=\operatorname{Sup}(\operatorname{card}}$}
$\pi'_0=\operatorname{Sup}(\{\operatorname{card}A,(\unsure{\pi'_\alpha})_{\alpha\in A}\})$,
et soit $\pi\geq\unsure{2^{{\pi'_0}^{\pi_0}}}$ \deleted{\unsure{infini tel que} $\unsure{\pi\geq}$}
un cardinal tel que $\pi^{\pi'_0}=\pi$. Alors $S$
est $\pi$-accessible, i.e. $\underline S$ est $\pi$-accessible dans $M$.
On aura en effet} Soit
\[
\begin{aligned}
\pi'_0&=\deleted{\unsure{\operatorname{Sup}(\operatorname{card}A,}}\operatorname{Sup}_{\alpha\in A}\pi'_\alpha \deleted{\unsure{,\operatorname{card} A})}\\
\pi''_0&=\operatorname{Sup}(\operatorname{card}A\deleted{\unsure{,\pi}},\operatorname{Sup}_\alpha 2^{{\pi'_\alpha}^{\pi_0}})
\end{aligned}
\]
% slanted
\margin{Ça a l'air un peu \unsure{abracadabrant} ! Il faut de plus supposer que $\pi$ n'est pas cardinal \unsure{et|$\infty$|$\aleph$} limite !}
et soit $\pi\geq\pi''_0$ un cardinal tel que
$\pi^{\pi'_0}=\pi$, i.e. de la forme $c^{\pi'_0}$, avec $c\geq\pi'_0$.
Alors $S$ est $\pi$-accessible (i.e. $\underline S$ est $\pi$-accessible
dans $M$). En particulier, $S$ est $2^{\pi''_0}$-accessible.
